\section{Related Work}
\label{sec:related}

\paragraph{Learned multi-view geometry.} Feed-forward models for multi-view 3D reconstruction have advanced rapidly. DUSt3R~\cite{dust3r} introduced pairwise pointmap regression; MASt3R~\cite{mast3r} extended it with local feature matching. Pi3X~\cite{pi3x} and VGGT~\cite{vggt} process all $N$ views jointly via permutation-equivariant architectures, achieving stronger global consistency. These models output camera poses and dense geometry from images alone, but their accuracy on wide-baseline, feature-adversarial scenes remains limited.

\paragraph{Multi-camera calibration.} Classical extrinsic calibration uses physical targets (checkerboards~\cite{zhang2000}), vanishing points~\cite{cipolla1999}, or structure-from-motion pipelines~\cite{schonberger2016}. Target-free methods assume sufficient static texture for feature matching across views---an assumption violated in dynamic, textureless environments. Our setting falls outside the scope of these methods.

\paragraph{Pose graph optimization.} Rotation averaging~\cite{hartley2013,chatterjee2018} and translation averaging~\cite{moulon2013} are classical tools for aggregating noisy pairwise relative poses into a globally consistent estimate. We adopt these as our post-processing backbone (Stage~1) and as the fusion mechanism for combining heterogeneous edge sources (Stages~3--4).

\paragraph{Test-time refinement.} Test-time training (TTT)~\cite{sun2020} and test-time adaptation adapt model predictions at inference time using self-supervised objectives. In geometry estimation, self-supervised depth refinement~\cite{godard2019} fine-tunes networks on test sequences using photometric consistency. Our approach differs in that we do \emph{not} modify model weights---instead, we exploit the model's conditioning interface to iteratively refine predictions, treating the pretrained model as a fixed denoiser.

\paragraph{Iterative pose refinement.} Classical approaches such as ICP~\cite{besl1992} and PnP+RANSAC refinement iterate between correspondence estimation and pose optimization. Bundle adjustment~\cite{triggs2000} jointly optimises poses and 3D structure. Our method is conceptually similar---iterative refinement toward a fixed point---but operates through a learned model's conditioning mechanism rather than explicit geometric optimization.

\paragraph{Denoising and diffusion on manifolds.} Score-based diffusion models~\cite{song2021} learn to denoise data via iterative refinement from noise toward the data manifold. Recent work has extended diffusion to SE(3)~\cite{yim2023} for molecular and pose generation. Our Monte Carlo perturbation strategy (Section~\ref{sec:stage4}) shares this structure: perturb the current estimate, apply a learned denoiser (Pi3X conditioning), and average. However, we operate in a discriminative setting (fixed pretrained model, no diffusion training) and exploit the observation that convergence is restricted to the SO(3) component.
